%! Author = lili
%! Date = 25.09.25

Additionally to the iGEM Slack, platforms such as \href{https://stackoverflow.com/}{Stackoverflow} or \href{https://www.reddit.com/}{Reddit} can be useful for finding answers to coding questions.

\section{Debugging}
Implementing various debugging strategies can help identify and resolve issues efficiently and reducing time spend
troubleshooting.

Recommended strategies:

\paragraph{Try-catch statements} are a classic method for handling errors in code.
This approach allows to attempt to execute a block of code and gracefully handle any exceptions that arise.
By capturing errors, developers can provide meaningful feedback or take corrective actions without crashing the application.
    \begin{mybox}
        \begin{lstlisting}[language=TypeScript]
try {
    {/* Code that may throw an error */}
    const response = await fetch(url);
    const data = await response.json();
} catch (error) {
    console.error("Error fetching data:", error);
    alert("There was an issue retrieving data. Please try again later.");
}
        \end{lstlisting}
    \end{mybox}

\paragraph{Console logging} is another straightforward yet effective debugging tool.
By placing \texttt{console.log()} statements throughout the code, you can track variable values and
application flows.
This practice is especially useful during the initial development of new components or when troubleshooting
specific issues, as it provides real-time feedback on how the application is behaving.
    \begin{mybox}
        \begin{lstlisting}[language=TypeScript]
console.log("Current state:", currentState);
        \end{lstlisting}
    \end{mybox}
\paragraph{Visual error detection} is not a classic approach, but it can be useful to catch errors that are not
syntax-related but arise from incorrect data or missing values.
For example, in situations where a form field must not be empty, visual error reminders can alert users to issues before submission. %Further information can be found in appendix part \hyperlink{sec:14.0.4}{4}.
\paragraph{Implementing data validation checks} is crucial for controlling datasets and preventing semantic errors.
By validating inputs before processing them, you can catch potential issues early.
This can involve checking for
required fields, ensuring data types match
expectations, and validating formats (e.g., url or phone number formats).
That is why I used TypeScript for my template and would recommend using TypeScript overall.
    \begin{mybox}
        \begin{lstlisting}[language=TypeScript]
function validateInput(data) {
    if (!data.email.includes("@")) {
        throw new Error("Invalid email format.");
    }
}
        \end{lstlisting}
    \end{mybox}

